\scanpage{77}
% archive foliation: 79
\diagram{\begin{tikzcd}[ampersand replacement=\&,column sep=small,row sep=large]
MDB_{g,\nu} \arrow[rr,dash] \arrow[dr] \&\& M_{\mathrm{st}}ND_{g,\nu} \arrow[dl] \\
\& mdt_{g,\nu} \&
\end{tikzcd}}
(dans $mdt$, la lettre $t$ rappelle qu'on a inclus
des structures toriques dans le type
combinatoire). Sûrement aussi, si on
\deleted{définit}\inserted{considère} $MD_{g,\nu}$ (sans bord, ou à bord
``généralisé'' réduit à un ens de $\nu$ pts)
et $M_{\mathrm{st}}N\deleted{D}_{g,\nu}$ (sans structure \inserted{multi}disque) on
aura
\diagram{\begin{tikzcd}[ampersand replacement=\&,column sep=small,row sep=large]
MD_{g,\nu} \arrow[rr] \arrow[dr] \&\& M_{\mathrm{st}}N_{g,\nu} \arrow[dl] \\
\& md_{g,\nu} \&
\end{tikzcd}}
cette fois un morphisme de \deleted{\unsure{de}}
multiplicités \deleted{\unsure{analyti.}} $\mathbb C$-algébriques.
Dans le cas 1), les dimensions \inserted{(réelles)} de
part et d'autre sont celles de $\mathbb C^{3g-3+\nu}\times\mathbb R^\nu$,
dans le cas 2), \inserted{les dim. complexes sont} celle de $\mathbb C^{3g-3+\nu}$.
